\chapter{Selenium}
    \section{Why Selenium?}
    \textbf{What are the benefits compared to the actual scripts? 
    Would there be other solutions/frameworks? Etc...}

    \section{Setting up Selenium}
    As Selenium is designed to work and test within web browsers, our setup 
    needs some extra components.\\
    Thankfully we already have a solution in the form of KDE's "at-spi" 
    accessability api.

    \subsection{Overwiew}
    \textbf{Selenium} A Browser automation framework.\\
    \textbf{Driver} The server component of Selenium.\\
    \textbf{Appium} An extension to Selenium making it more useful for app 
    testing.\\
    \textbf{AT-SPI2} A protocol over DBus, which GUI toolkit widgets use to 
    provide their content to screen readers.\\

    \subsection{Installing Dependencies}
    Some or most of the following packages might already been installed on 
    the fedora distro, but to make sure, we will run the install commands 
    through our package manger.
    \begin{itemize}
        \item a supported Python3 version
        \item at-spi2
        \item pyatspi2
        \item accerciser
        \item gobject-introspection
    \end{itemize}
    We are gonna install these and other dependencies through our package 
    manager ("DNF" on fedora):
    \begin{lstlisting}[language=bash, gobble=8]
        sudo dnf install \
        at-spi2-core \
        accerciser \
        gobject-introspection \
        python3 \
        python3-pip \
        python3-pyatspi \
        python3-gobject \
        at-spi2-core \
        accerciser \
        cairo-devel
    \end{lstlisting}
    Additionally we ar going to need to install the following packages:
    \begin{lstlisting}[language=bash, gobble=8]
        sudo dnf install gcc \
        gobject-introspection-devel \
        cairo-gobject-devel \
        pkg-config python3-devel \
        gtk4
    \end{lstlisting}
    Including the following in order to be able to build 
    "selenium-webdriver-at-spi" in the following steps.
    \begin{lstlisting}[language=bash, gobble=8]
        sudo dnf install \
        extra-cmake-modules \
        qt6-qtbase-devel \
        qt6-qtbase-private-devel \
        qt6-qtdeclarative-devel \
        kf6-kwindowsystem-devel \
        plasma-wayland-protocols-devel \
        kwayland-devel \
        kpipewire-devel \
        kwin-devel \
        wayland-devel \
        ruby
    \end{lstlisting}
    Once we have these packages installed we are now ready to clone and 
    install selenium.

    \subsection{Installing Selenium-webdriver-AT-SPI}
    \begin{lstlisting}[language=bash, gobble=8]
        git clone https://invent.kde.org/sdk/selenium-webdriver-at-spi.git
        cd selenium-webdriver-at-spi
        # For virtual environment
        python3 -m venv --system-site-packages venv
        source venv/bin/activate
        pip3 install -r requirements.txt
        echo export PATH="~/.local/bin:$PATH" >> ${HOME}/.bashrc
        mkdir build
        cd build
        cmake ..
        make
        sudo make install
        cd ..
    \end{lstlisting}
    Now that Selenium is built, it can be used to run the Selenium-specific test.

    \subsection{Running Selenium Tests}
    If we are running our tests inside a VM -- which we do when working 
    with the scripts locally --, we make sure to set it up properly.\\
    In this case we are using QEMU/KVM Virtual Machine Manager.\\

    In the VM's settings, we go to the video tab and make sure 
    the Mode is set to Virtio and that 3D acceleration is checked.\\
    Then under Display Spice we need to make sure OpenGL is checked and 
    listen type is set to None.\\

    Now that the VM is set up and after rebooting the guest machine, we 
    will be able to run the Selenium tests.
    In order to test if it works we can run the example test by using the 
    following command from the selenium folder.

    \begin{lstlisting}[language=bash, gobble=8]
        selenium-webdriver-at-spi-run ./examples/calculatortest.py
    \end{lstlisting}
    If everything is working correctly, a window will open in which the 
    KCalc app will open and the test will run.

    \subsection{Using already existing Selenium scripts}
    Given an KDE-repository containing Selenium tests scripts for KDE Linux with OpenQA, it will be checked, if these scripts are of any further use for the current project.
    The repository https://invent.kde.org/kde-linux/os-autoinst-distri-kdelinux/-/tree/master/tests has the following structure:
    \begin{lstlisting}[language=bash, gobble=8]
        tests/
        common/
        kdelinux/
        kdelinux-live/
    \end{lstlisting}
    Each folder contains tests for different cases. The repository defines three test suites:
    \begin{itemize}
        \item install-system
        \item upgrade-system
        \item sanity-test
    \end{itemize}
    Each test-suite contains multiple tests, but most of them are for the os, like for example booting the system.
    Though in 'kdelinux/app/' there are some scripts, which could be useful for this work.
    First there is 'firefox.py':

    \begin{itemize}
        \item shebang, chmod u+x firefox.py
        \item create venv, install unittest-xml-reporting
        \item copy lib folder to where the script lies (!!find workaround!!)
    \end{itemize}

    Follow steps from \#\# Run tests locally - install podman-compose
